On dispose d'une solution d'acide bromhydrique telle que
$[\ce{H3O+}]=\qty{2.5e-3}{\mol\per\L}$.
\begin{enumerate}
    \item Écrire l'équation de la réaction de dissolution du bromure d'hydrogène
          \ce{HBr(g)} dans l'eau.
    \item En déduire la concentration $[\ce{Br-(aq)}]$ et la concentration
          molaire de la solution.
    \item Quel volume d'eau faut-il ajouter à un volume $V=\qty{50}{\mL}$ de
          cette solution pour obtenir une solution dans laquelle
          $[\ce{H3O+}]=\qty{5e-4}{\mol\per\L}$~?
\end{enumerate}
\textbf{Indication~:} Le bromure d'hydrogène \ce{HBr} est un gaz soluble dans
l'eau~; après dissolution, sa solution aqueuse constitue l'acide brombydrique.

